\documentclass[10.5pt]{config}
% 本实验报告采用 署名-相同方式共享 4.0 国际许可协议 (CC BY-SA 4.0)
\usepackage{multirow}
\usepackage{wrapfig}
\usepackage{setspace}
\usepackage{titlesec}
\usepackage{ctex}
\usepackage{graphicx} % 用于插入图片
\usepackage{float}    % 用于控制图片位置
\usepackage{makecell}
\usepackage{tikz}
\usepackage{enumitem} % 用于更灵活的列表环境
\usetikzlibrary{arrows.meta}

\titleformat{\subsubsection}[block]
{\normalfont\bfseries}{\thesubsubsection}{1em}{}
\titlespacing*{\subsubsection}{1.5em}{3.25ex}{1.5ex}
\setstretch{1.5}

% --- 请在这里填写您的个人信息 ---
\major{机械工程} % 专业
\name{韩颂义}
\partner{} % 签名栏
\stuid{3240104198}
\date{5.28} % 可根据实际实验日期修改，如 2026.5.22
\lab{东3-308}
\grades{}
\course{电工电子学实验}
\instructor{季瑞松} % 指导老师
\exptype{基础实验} % 实验类型
\expname{直流稳压电源} % 按照本次实验要求修改实验名称

\begin{document}

\makeheader % 首页头部

\section{实验目的}

\begin{enumerate}[label=\arabic*., leftmargin=2em]
    \item 掌握单相桥式整流电路的工作原理。
    \item 观察几种常用滤波电路的效果。
    \item 理解集成稳压器的工作原理和使用方法。
    \item 掌握直流稳压电源主要技术指标的测试方法。
\end{enumerate}

\section{实验原理}

\subsection*{1. 单相桥式整流电路}
单相桥式整流电路由四个二极管组成桥式结构，使交流电压的正、负半周均通过负载形成同一方向的电流，从而得到脉动直流电压。理想情况下，桥式整流电路输出电压平均值 $U_L$ 与变压器副边交流电压有效值 $U_2$ 近似满足：
$$U_L \approx 0.9U_2$$

\subsection*{2. 滤波电路}
桥式整流后得到的直流电压仍含有较大的脉动分量，可利用电容的储能作用使输出电压平滑化。滤波电容越大，放电时间越长，输出纹波通常越小。桥式整流加电容滤波时，负载上的电压平均值 $U_L$ 与变压器副边电压有效值 $U_2$ 的关系可近似表示为：
$$U_L \approx 1.2U_2$$

\subsection*{3. 稳压电路}
经过滤波得到的直流电压会随着交流电源电压或负载电流变化而变化。为了获得稳定的直流输出电压，需要在整流滤波电路后加入稳压电路。本实验使用三端集成稳压器 7812，其三个引脚分别为输入端、接地端和输出端。使用时应保证输入端电压高于输出端电压，并留有足够的输入输出压差。

\subsection*{4. 主要性能指标}
\begin{enumerate}[label=(\arabic*), leftmargin=2em]
    \item \textbf{纹波系数 $\gamma$} \\
    用来表征电路输出电压的脉动程度，定义为输出电压中交流分量有效值 $\tilde{U}_L$ 与输出电压平均值 $U_L$ 之比，即 $\gamma = \frac{\tilde{U}_L}{U_L}$。$\gamma$ 值越小，输出电压越平滑。
    \item \textbf{输出电压 $U_o$ 和输出电流 $I_o$} \\
    $U_o$ 通常指稳压后的额定直流输出电压值，$I_o$ 通常指稳压器的额定输出电流值。
    \item \textbf{输出电阻 $r_o$} \\
    当输入交流电压 $U_2$ 保持不变时，输出电阻可由输出电压变化量与输出电流变化量之比求得，即 $r_o = \frac{\Delta U_L}{\Delta I_L}$。
    \item \textbf{稳压系数 $S$} \\
    当负载保持不变时，稳压系数表示输出电压相对变化量与输入电压相对变化量之比，即 $S = \frac{\Delta U_L / U_L}{\Delta U_I / U_I}$。$S$ 越小，稳压性能越好。
\end{enumerate}

\section{实验内容与数据记录}

\subsection*{1. 单相整流、滤波电路}
取变压器副边电压 15V 挡作为整流电路的输入电压 $u_2$，并实测 $U_2$ 的值。完成下表的测量。

\begin{table}[H]
    \centering
    \caption{单相整流、滤波电路测试记录 ($R_L = 240\Omega$, $U_2 = 16.1375$ V)}
    \begin{tabular}{|c|c|c|c|}
        \hline
        \multirow{2}{*}{电路图说明} & \multicolumn{2}{c|}{测量结果} & 计算值 \\ \cline{2-4}
        & $U_L$ / V & $\tilde{U}_L$ / V & $\gamma$ \\ \hline
        桥式整流（无滤波） & 13.3630 & 7.2499 & 0.5425 \\ \hline
        半波整流 + $100\mu\text{F}$ 滤波 & 19.4151 & 1.56017 & 0.0804 \\ \hline
        半波整流 + $470\mu\text{F}$ 滤波 & 20.079 & 0.38192 & 0.0190 \\ \hline
        桥式整流 + $\pi$型RC滤波 & 22.497 & 0.27055 & 0.0120 \\ \hline
    \end{tabular}
\end{table}

% --- 插入单相整流、滤波电路的波形图 ---
\begin{figure}[H]
    \centering
    \includegraphics[width=0.5\textwidth]{1.jpg}
    \caption{桥式整流（无滤波）输出波形}
    \label{fig:wave_rect_1}
\end{figure}

\begin{figure}[H]
    \centering
    \includegraphics[width=0.5\textwidth]{2.jpg}
    \caption{半波整流加 $100\mu\text{F}$ 滤波输出波形}
    \label{fig:wave_rect_2}
\end{figure}

\begin{figure}[H]
    \centering
    \includegraphics[width=0.5\textwidth]{3.jpg}
    \caption{半波整流加 $470\mu\text{F}$ 滤波输出波形}
    \label{fig:wave_rect_3}
\end{figure}

\begin{figure}[H]
    \centering
    \includegraphics[width=0.5\textwidth]{4.jpg}
    \caption{桥式整流加 $\pi$ 型 RC 滤波输出波形}
    \label{fig:wave_rect_4}
\end{figure}

\subsection*{2. 集成稳压电路}
\subsubsection*{(1) 改变负载 $R_L$ 测试输出电阻 $r_o$}
取变压器副边电压 15V 挡作为输入电压 $u_2$，按图接好整流、滤波、稳压(7812)电路，改变负载电阻 $R_L$。

\begin{table}[H]
    \centering
    \caption{集成稳压电路负载变化测试 ($U_2 = 11.8763$ V)}
    \begin{tabular}{|c|c|c|c|c|}
        \hline
        负载 & \multicolumn{3}{c|}{测量结果} & 计算值 \\ \hline
        $R_L / \Omega$ & $U_L$ / V & $\tilde{U}_L$ / mV & $I_L$ / mA & $r_o$ \\ \hline
        $\infty$ & 11.7754 & 0.217 & 0 & x \\ \hline
        240 & 11.7535 & 0.194 & 49.2 & $0.445\,\Omega$ \\ \hline
        120 & 11.7394 & 0.304 & 98.5 & $0.365\,\Omega$ \\ \hline
    \end{tabular}
\end{table}

% --- 插入负载变化时的波形图 ---
\begin{figure}[H]
    \centering
    \includegraphics[width=0.5\textwidth]{11.jpg}
    \caption{空载时稳压电路输出波形}
    \label{fig:wave_load_1}
\end{figure}

\begin{figure}[H]
    \centering
    \includegraphics[width=0.5\textwidth]{22.jpg}
    \caption{$R_L=240\Omega$ 时稳压电路输出波形}
    \label{fig:wave_load_2}
\end{figure}

\begin{figure}[H]
    \centering
    \includegraphics[width=0.5\textwidth]{33.jpg}
    \caption{$R_L=120\Omega$ 时稳压电路输出波形}
    \label{fig:wave_load_3}
\end{figure}

\subsubsection*{(2) 改变输入电压 $u_2$ 测试稳压系数 $S$}
取负载电阻 $R_L = 120\Omega$ 不变，改变电路输入电压 $u_2$ 的抽头。

\begin{table}[H]
    \centering
    \caption{集成稳压电路输入变化测试 ($R_L = 120\Omega$)}
    \begin{tabular}{|c|c|c|c|c|}
        \hline
        \multirow{2}{*}{变压器抽头} & \multicolumn{3}{c|}{测量结果} & 计算值 \\ \cline{2-5}
        & $U_2$ / V & $U_L$ / V & $\tilde{U}_L$ / mV & $S$ \\ \hline
        9V 挡 & 9.8772 & 11.7931 & 404.1 &  \\ \hline
        12V 挡 & 13.0632 & 11.7892 & 0.224 & $0.0009$ \\ \hline
        15V 挡 & 16.2842 & 11.7921 & 0.304 & $0.0004$ \\ \hline
        18V 挡 & 19.5756 & 11.7924 & 0.209 & $0.0003$ \\ \hline
    \end{tabular}
\end{table}

% --- 插入输入电压变化时的波形图 ---
\begin{figure}[H]
    \centering
    \includegraphics[width=0.6\textwidth]{111.jpg}
    \caption{9V 挡时稳压电路输出波形}
    \label{fig:wave_input_1}
\end{figure}

\begin{figure}[H]
    \centering
    \includegraphics[width=0.6\textwidth]{222.jpg}
    \caption{12V 挡时稳压电路输出波形}
    \label{fig:wave_input_2}
\end{figure}

\begin{figure}[H]
    \centering
    \includegraphics[width=0.6\textwidth]{333.jpg}
    \caption{15V 挡时稳压电路输出波形}
    \label{fig:wave_input_3}
\end{figure}

\begin{figure}[H]
    \centering
    \includegraphics[width=0.6\textwidth]{444.jpg}
    \caption{18V 挡时稳压电路输出波形}
    \label{fig:wave_input_4}
\end{figure}

\section{实验总结}

\begin{enumerate}[label=\arabic*., leftmargin=2em]
    \item \textbf{根据表 5.22-1 的结果，讨论桥式整流电路输出电压平均值 $U_L$ 和输入交流电压有效值 $U_2$ 之间的数值关系。}
    \\ \textbf{解答：}
    理想情况下，桥式整流电路在不加滤波电容时，输出电压平均值与变压器副边交流有效值之间近似满足 $U_L \approx 0.9U_2$。本实验中 $U_2 = 16.1375\text{V}$，按理想关系估算应有 $U_L \approx 14.52\text{V}$，而表中桥式整流无滤波时测得 $U_L = 13.3630\text{V}$，实测值与理论估算基本接近但略小。产生偏差的原因可能包括二极管导通压降、变压器输出波动、测量仪表读数方式以及实验接线和波形读数误差等。总体来看，桥式整流能把交流电压变换为单向脉动直流电压，其输出平均值与输入交流有效值之间符合近似比例关系。

    \item \textbf{根据表 5.22-1 的结果，总结不同滤波电路的滤波效果。}
    \\ \textbf{解答：}
    由表 5.22-1 可知，无滤波桥式整流电路的交流纹波分量较大，纹波系数为 $0.5425$，输出电压脉动明显。加入 $100\mu\text{F}$ 滤波电容后，纹波系数降为 $0.0804$，说明电容能够在电压下降阶段向负载放电，从而减小输出波动；当电容增大到 $470\mu\text{F}$ 时，纹波系数进一步降为 $0.0190$，滤波效果明显增强。桥式整流加 $\pi$ 型 RC 滤波时，纹波系数最小，为 $0.0120$，说明多级滤波对纹波抑制更充分，输出电压更接近平滑直流。因此，在本实验条件下，滤波效果由弱到强大致为：无滤波、$100\mu\text{F}$ 电容滤波、$470\mu\text{F}$ 电容滤波、$\pi$ 型 RC 滤波。

    \item \textbf{根据表 5.22-2 和表 5.22-3 结果，分析集成稳压器的稳压性能。}
    \\ \textbf{解答：}
    从负载变化测试看，负载由空载变为 $240\Omega$ 和 $120\Omega$ 时，输出电压分别为 $11.7754\text{V}$、$11.7535\text{V}$ 和 $11.7394\text{V}$，变化量较小，计算得到输出电阻分别约为 $0.445\Omega$ 和 $0.365\Omega$，说明 7812 集成稳压器在负载电流变化时仍能使输出电压保持在约 $11.7\text{V}$，具有较好的负载调整能力。从输入电压变化测试看，输入电压在 12V、15V、18V 挡变化时，输出电压基本稳定在 $11.79\text{V}$ 附近，稳压系数数量级较小，说明稳压器能够有效抑制输入电压变化。但 9V 挡时纹波电压达到 $404.1\text{mV}$，明显大于其他挡位，说明输入电压较低时稳压器输入输出压差不足，稳压性能会下降。因此，实际使用三端稳压器时，应保证输入端电压高于输出电压并留有足够压差余量。

    \item \textbf{总结实验中出现的问题及其解决方法。}
    \\ \textbf{解答：}
    本次实验中容易出现的问题主要有以下几类：首先，整流、滤波和稳压电路连接环节较多，二极管方向、电解电容极性若接错，会导致输出电压异常，甚至可能损坏元件；其次，纹波电压通常较小，如果示波器耦合方式、垂直灵敏度或时基设置不合适，容易造成波形显示不清或读数误差；另外，多次更换负载和输入抽头时，若未及时记录对应条件，也容易造成数据混淆。针对这些问题，实验时应先在断电状态下检查电路连接和元件极性，确认无误后再通电测量；测量直流输出时可先用万用表判断电压范围，再用示波器观察纹波波形；读取纹波时应根据需要选择合适的耦合方式和量程，并保持每组数据的输入电压、负载和波形记录一一对应。通过本实验，可以更直观地理解整流、滤波、稳压各环节的作用，也体会到规范操作和细致记录对实验结果可靠性的影响。
\end{enumerate}

\end{document}